% registry: EK
\methodhead{EK}{Ellyin and Kujawski (1993): multiaxial criterion with mean-stress effect}\label{sheet:EK}

\block{Source} F.~Ellyin, D.~Kujawski, in \emph{Advances in Multiaxial Fatigue}, ASTM STP 1191
(1993) 55--66 \cite{ellyin1993}, Eqs.~(5b), (9)--(18), pp.~60--61; Fig.~2.

\block{Idea and scope} The damage parameter is the cyclic distortional strain-energy range, raised
by a mean-stress factor built on the sum of the principal mean stresses and divided by a
multiaxial constraint factor computed from the free-surface principal strains. The paper applies
the high-cycle form (Eqs.~16--18) ``when bulk material response is essentially elastic''; it does
not treat non-proportional cycles.

\block{Inputs}
\begin{center}\small
\begin{tabular}{@{}llp{100mm}@{}}
\toprule
Symbol & Unit & Meaning\\\midrule
$\Delta W_d$, $\Delta W_d^e$, $\Delta W_d^p$ & MPa & distortional strain-energy range per reversal; elastic and plastic parts\\
$\bar\sigma$, $\Delta\bar\sigma=2\bar\sigma$ & MPa & effective (von Mises) stress amplitude and range\\
$\sigma_m$ & MPa & sum of the principal mean stresses, $\sigma_{1,m}+\sigma_{2,m}+\sigma_{3,m}=\tr\Sm$ (Eq.~5b)\\
$\sigma_a$, $\sigma_{-1}$ & MPa & uniaxial stress amplitude and fully reversed uniaxial strength (Eq.~18)\\
$E$, $\nu=\nu_e$ & MPa, -- & Young's modulus and elastic Poisson's ratio\\
$E^*$ & MPa & $3E/[2(1+\nu)]$\\
$\bar\rho$ & -- & multiaxial constraint factor, Eq.~\eqref{eq:ek-rho}\\
$\bar\nu$, $\nu_p$ & -- & effective Poisson's ratio (Eq.~13) and plastic Poisson's ratio (0.5)\\
$\varepsilon_{\max}$, $\gamma_{\max}$ & -- & maximum in-plane principal strain and maximum shear strain over the cycle\\
$\varepsilon_r$ & -- & normal strain of the free surface (Eq.~14)\\
$\varepsilon_{xx}$, $\varepsilon_{yy}$, $\varepsilon_{zz}$ & -- & surface strains ($x$--$y$ the surface, $z$ its normal)\\
$\alpha$, $n$ & -- & mean-stress weight (\texttt{alpha\_EK}) and constraint exponent (\texttt{n\_EK})\\
$K'$, $n'$ & MPa, -- & optional cyclic Ramberg--Osgood coefficient and exponent (\texttt{K\_prime}, \texttt{n\_prime})\\
$\Delta\varepsilon_p$, $\varepsilon_{p,a}$ & -- & plastic strain range and amplitude\\
$\psi$ & MPa & damage parameter (MJ/m$^3$)\\
\bottomrule
\end{tabular}
\end{center}

\block{Formulation} Ellyin and Kujawski (Eqs.~11--18) define
\begin{equation}
\psi = \frac{\Delta W_d}{\bar\rho^{\,n}}\Bigl(1+\alpha\,\frac{\sigma_m}{\bar\sigma}\Bigr),\qquad
\Delta W_d^e = \frac{\Delta\bar\sigma^2}{2E^*},\qquad E^* = \frac{3E}{2(1+\nu)} ,
\label{eq:ek-main}
\end{equation}
with $\sigma_m$ the sum of the principal mean stresses, without a factor $1/3$ (Eq.~5b), and the
constraint factor (Eq.~12)
\begin{equation}
\bar\rho = (1+\bar\nu)\,\frac{\varepsilon_{\max}}{\gamma_{\max}} ,
\label{eq:ek-rho}
\end{equation}
where $\bar\nu=\nu_e$ for elastic strains (Eq.~13). Its limiting values are $\bar\rho=1$ uniaxial,
$1+\nu$ in reversed torsion and $1-\nu$ equibiaxial (p.~60). The paper sets $\alpha=1/2$ and $n=1$
(p.~61, Eq.~18); in uniaxial loading this gives
$\sigma_a/\sigma_{-1}=\bigl(1+\sigma_m/(2\sigma_a)\bigr)^{-1/2}$. The repository computes
$\bar\sigma=\sqrt3\sqrt{J_{2,a}}$, $\varepsilon_{\max}$ and $\gamma_{\max}$ from the surface strains
with $\varepsilon_{zz}=-\nu(\varepsilon_{xx}+\varepsilon_{yy})/(1-\nu)$ (Eq.~14), the maxima taken
over 360 phase samples of the cycle.

This is not a critical-plane method: no plane search, tie rule or \texttt{np\_shear\_rule}
applies, and the global shear-amplitude selection does not change it.

\block{Constants and calibration} $\alpha=0.5$ and $n=1$ by default, as in the paper; both can be
supplied (\texttt{alpha\_EK}, \texttt{n\_EK}). $K'$, $n'$ are optional and never fitted. No EK
constant is calibrated statistically; only the master curve is fitted.

\block{Life} The repository uses the master curve on $\psi$, Eq.~\eqref{eq:mastercurve}, fitted on
the uniaxial $R=-1$ design set.

\block{Repository choices}
\begin{itemize}
\item High-cycle elastic form by default, $\Delta W_d=\Delta W_d^e$: $K'$, $n'$ are rarely present in
  HCF workbooks and never fitted. When both are supplied, the Masing plastic term
  $\Delta W_d^p=\frac{1-n'}{2(1+n')}\,\Delta\bar\sigma\,\Delta\varepsilon_p$ with
  $\varepsilon_{p,a}=(\bar\sigma/K')^{1/n'}$, $\nu_p=0.5$ and a strain-weighted $\bar\nu$ (equal to
  Eq.~13 in uniaxial tension) is added under assumed stress control.
\item Free surface: $x$--$y$ is the surface plane and $\varepsilon_{zz}$ follows from plane stress.
\item Non-proportional single harmonics (not treated in the paper): $\bar\sigma$ is $\sqrt3$ times
  the longest deviatoric half-chord, and $\varepsilon_{\max}$, $\gamma_{\max}$ are maxima over the
  period.
\item Mean stress: Eq.~(18) is stated for $-1.06\le\sigma_m/\sigma_a\le4.4$; a non-positive factor
  $1+\alpha\sigma_m/\bar\sigma$ sets $\psi=0$.
\end{itemize}

\block{Worked example} Steel, $E=206\,000\MPa$, $\nu=0.3$, zero mean: uniaxial $300\MPa$ gives
$\bar\rho=1.00$, $\psi=0.75728\MPa$; torsion $200\MPa$ gives $\bar\rho=1.30$, $\psi=0.77670\MPa$.

\block{Verification} \texttt{tests/methods/test\_ellyin\_kujawski.py} (Eqs.~(5b), (9)--(18) and
Fig.~2 checked on the paper): $\bar\rho=1$, $1+\nu$, $1-\nu$ for uniaxial, torsional and equibiaxial states; Eq.~(13)
gives $\nu_e$ for elastic strains; the Eq.~(18) uniaxial mean-stress curve from
$\sigma_m/\sigma_a=-0.8$ to $4.0$; the torsion-to-tension ratio $\sqrt{(1+\nu)/3}$ ($n=1$) and
$1/\sqrt3$ ($n=0$); no effect of a mean shear stress; phase-origin invariance; the Masing term
recovers Eq.~(13) and the uniaxial closed-form energy. The paper prints no per-test values.
